\documentclass[10pt,a4paper]{amsart}
\usepackage{amsmath,amssymb,amsthm,graphicx}
\usepackage{verbatim}
\newenvironment{reference}{\comment}{\endcomment}
%\newenvironment{reference}{}{}
\newenvironment{slogan}{\comment}{\endcomment}
\newenvironment{history}{\comment}{\endcomment}

% For commutative diagrams we use Xy-pic
\usepackage[all]{xy}

% We use 2cell for 2-commutative diagrams.
\xyoption{2cell}
% Theorem environments.
%
\theoremstyle{plain}
\newtheorem{theorem}[subsection]{Theorem}
\newtheorem{proposition}[subsection]{Proposition}
\newtheorem{lemma}[subsection]{Lemma}

\theoremstyle{definition}
\newtheorem{definition}[subsection]{Definition}
\newtheorem{example}[subsection]{Example}
\newtheorem{exercise}[subsection]{Exercise}
\newtheorem{situation}[subsection]{Situation}

\theoremstyle{remark}
\newtheorem{remark}[subsection]{Remark}
\newtheorem{remarks}[subsection]{Remarks}

\numberwithin{equation}{subsection}

% Macros
%
\def\lim{\mathop{\mathrm{lim}}\nolimits}
\def\colim{\mathop{\mathrm{colim}}\nolimits}
\def\Spec{\mathop{\mathrm{Spec}}}
\def\Hom{\mathop{\mathrm{Hom}}\nolimits}
\def\Ext{\mathop{\mathrm{Ext}}\nolimits}
\def\SheafHom{\mathop{\mathcal{H}\!\mathit{om}}\nolimits}
\def\SheafExt{\mathop{\mathcal{E}\!\mathit{xt}}\nolimits}
\def\Sch{\mathit{Sch}}
\def\Mor{\mathop{\mathrm{Mor}}\nolimits}
\def\Ob{\mathop{\mathrm{Ob}}\nolimits}
\def\Sh{\mathop{\mathit{Sh}}\nolimits}
\def\NL{\mathop{N\!L}\nolimits}
\def\CH{\mathop{\mathrm{CH}}\nolimits}
\def\proetale{{pro\text{-}\acute{e}tale}}
\def\etale{{\acute{e}tale}}
\def\QCoh{\mathit{QCoh}}
\def\Ker{\mathop{\mathrm{Ker}}}
\def\Im{\mathop{\mathrm{Im}}}
\def\Coker{\mathop{\mathrm{Coker}}}
\def\Coim{\mathop{\mathrm{Coim}}}

% Boxtimes
%
\DeclareMathSymbol{\boxtimes}{\mathbin}{AMSa}{"02}

%
% Macros for moduli stacks/spaces
%
\def\QCohstack{\mathcal{QC}\!\mathit{oh}}
\def\Cohstack{\mathcal{C}\!\mathit{oh}}
\def\Spacesstack{\mathcal{S}\!\mathit{paces}}
\def\Quotfunctor{\mathrm{Quot}}
\def\Hilbfunctor{\mathrm{Hilb}}
\def\Curvesstack{\mathcal{C}\!\mathit{urves}}
\def\Polarizedstack{\mathcal{P}\!\mathit{olarized}}
\def\Complexesstack{\mathcal{C}\!\mathit{omplexes}}
% \Pic is the operator that assigns to X its picard group, usage \Pic(X)
% \Picardstack_{X/B} denotes the Picard stack of X over B
% \Picardfunctor_{X/B} denotes the Picard functor of X over B
\def\Pic{\mathop{\mathrm{Pic}}\nolimits}
\def\Picardstack{\mathcal{P}\!\mathit{ic}}
\def\Picardfunctor{\mathrm{Pic}}
\def\Deformationcategory{\mathcal{D}\!\mathit{ef}}
\usepackage{cleveref}
\crefname{originallemma}{original lemma}{original lemmas}
\Crefname{originallemma}{Original lemma}{Original lemmas}
\makeatletter
\newlabel{lemma-quasi-compact}{{00E8}{source}{}{sourcepointer-1}{}}
\newlabel{lemma-quasi-compact@cref}{{[originallemma][0][]00E8}{[1][0][]source}{}{}{}}
\makeatother
\crefname{originallemma}{original lemma}{original lemmas}
\Crefname{originallemma}{Original lemma}{Original lemmas}
\makeatletter
\newlabel{lemma-standard-etale-finite-flat-Zariski}{{00UF}{source}{}{sourcepointer-2}{}}
\newlabel{lemma-standard-etale-finite-flat-Zariski@cref}{{[originallemma][0][]00UF}{[1][0][]source}{}{}{}}
\makeatother
\crefname{originallemma}{original lemma}{original lemmas}
\Crefname{originallemma}{Original lemma}{Original lemmas}
\makeatletter
\newlabel{proposition-etale-locally-standard}{{00UE}{source}{}{sourcepointer-3}{}}
\newlabel{proposition-etale-locally-standard@cref}{{[originallemma][0][]00UE}{[1][0][]source}{}{}{}}
\makeatother
\title{A surjective etale map is covered by a finite flat one factoring through it locally}
\author{The Stacks Project Authors}
\date{stacks-project:algebra.tex:local-snapshot-20261012; GNU FDL 1.2 or later}
\begin{document}
\maketitle
\noindent\textit{Editorial note.} All delivered mathematical text is the author's own, copied from the pinned Stacks Project chapter file algebra.tex; no statement or proof step was rewritten, reordered or removed.\par\medskip
\noindent\textit{Editorial note.} Preamble provenance, as required for this source: the frozen snapshot stages the chapter's own preamble file beside it (preamble.tex, 217 lines, 8116 bytes, the same shared Stacks-wide preamble the published chow and etale-cohomology cards use); the first line of the pinned chapter file is \textbackslash{}input\{preamble\}, so no preamble text lives inside the chapter itself. The card takes that file's lines 8-18 (the verbatim package and the three \textbackslash{}comment-based environments reference/slogan/history at lines 9-12, section 13.2, together with the Xy-pic setup) and its lines 154 to the runtime end of the file (the native theorem declarations and mathematical macros). Not one character was changed, no package was added, no macro was defined, and the excerpt serves only as preamble, never as a span coordinate.\par\medskip
\noindent\textit{Editorial note.} The same-chapter pointers are located at their \textbackslash{}label lines inside the pinned chapter file; the cross-chapter pointers are located in the project's own tag list staged in the snapshot (sources/stacks-chow/tags.txt), at the verbatim lines carrying their canonical tags.\par\medskip
\noindent\textit{Editorial note.} The counted proof refers to same-chapter statements by label; they are recorded as navigation stubs pointed at their \textbackslash{}label lines, and the arguments using them stay inside the counted span.\par\medskip
\noindent\textit{Editorial note.} The counted spans contain no \textbackslash{}cite, so the card carries no bibliography.\par\medskip
\noindent\textit{Allowed standard prerequisites (editorial).}\begin{itemize}
\item Local standard etaleness over a finite cover. Used at the start of the counted proof (lines 40367-40370). Produces finitely many elements generating the unit ideal over which the map is standard etale.
\item Products of finite flat algebras. Used to assemble the global statement (lines 40370-40390). Combines the local solutions into one finite flat algebra mapping surjectively to the base.
\end{itemize}\medskip
\ifdefined\hypertarget\hypertarget{sourcepointer-1}{}\fi
\noindent\textit{Source reference.} Original lemma 00E8: sources/stacks-algebra/pinned-original-source.tex, source lines 3469--3469. Navigation pointer for \textbackslash{}ref\{lemma-quasi-compact\} used inside the counted span (line 40368); the argument that consumes it is retained in the span.\par\medskip
\ifdefined\hypertarget\hypertarget{sourcepointer-2}{}\fi
\noindent\textit{Source reference.} Original lemma 00UF: sources/stacks-algebra/pinned-original-source.tex, source lines 40287--40287. Navigation pointer for \textbackslash{}ref\{lemma-standard-etale-finite-flat-Zariski\} used inside the counted span (line 40378); the argument that consumes it is retained in the span.\par\medskip
\ifdefined\hypertarget\hypertarget{sourcepointer-3}{}\fi
\noindent\textit{Source reference.} Original lemma 00UE: sources/stacks-algebra/pinned-original-source.tex, source lines 40061--40061. Navigation pointer for \textbackslash{}ref\{proposition-etale-locally-standard\} used inside the counted span (line 40367); the argument that consumes it is retained in the span.\par\medskip
\begin{lemma}
\label{lemma-etale-finite-flat-zariski}
Let $R \to S$ be a ring map.
Assume that
\begin{enumerate}
\item $R \to S$ is \'etale, and
\item $\Spec(S) \to \Spec(R)$ is surjective.
\end{enumerate}
Then there exists a ring map $R \to S'$ such that
\begin{enumerate}
\item $R \to S'$ is finite, finitely presented, and flat
(in other words it is finite projective as an $R$-module),
\item $\Spec(S') \to \Spec(R)$ is surjective,
\item for every prime $\mathfrak q' \subset S'$ there exists a
$g' \in S'$, $g' \not \in \mathfrak q'$ such that
the ring map $R \to S'_{g'}$ factors as $R \to S \to S'_{g'}$.
\end{enumerate}
\end{lemma}

\begin{proof}
By Proposition \ref{proposition-etale-locally-standard} and
the quasi-compactness of $\Spec(S)$ (see Lemma \ref{lemma-quasi-compact})
we can find $g_1, \ldots, g_n \in S$ generating the unit ideal
of $S$ such that each $R \to S_{g_i}$ is standard \'etale.
If we prove the lemma for the ring map $R \to \prod_{i = 1, \ldots, n} S_{g_i}$
then the lemma follows for the ring map $R \to S$.
Hence we may assume that $S = \prod_{i = 1, \ldots, n} S_i$
is a finite product of standard \'etale morphisms.

\medskip\noindent
For each $i$ choose a ring map $R \to S_i'$ as in
Lemma \ref{lemma-standard-etale-finite-flat-Zariski}
adapted to the standard \'etale morphism $R \to S_i$.
Set $S' = S_1' \otimes_R \ldots \otimes_R S_n'$; we will use
the $R$-algebra maps $S_i' \to S'$ without further mention below.
We claim this works. Properties (1) and (2) are immediate.
For property (3) suppose that $\mathfrak q' \subset S'$ is a prime.
Denote $\mathfrak p$ its image in $\Spec(R)$.
Choose $i \in \{1, \ldots, n\}$ such that $\mathfrak p$
is in the image of $\Spec(S_i) \to \Spec(R)$; this is
possible by assumption. Set $\mathfrak q_i' \subset S_i'$
the image of $\mathfrak q'$ in the spectrum of $S_i'$.
By construction of $S'_i$ there exists a $g'_i \in S_i'$
such that $R \to (S_i')_{g_i'}$ factors as
$R \to S_i \to (S_i')_{g_i'}$. Hence also
$R \to S'_{g_i'}$ factors as
$$
R \to S_i \to (S_i')_{g_i'} \to S'_{g_i'}
$$
as desired.
\end{proof}

\end{document}
